\chapter{Convolución en el espacio euclidiano}\label{cap:v4-06}
Las estimaciones de convolución no dependen esencialmente de trabajar en una sola variable. Desarrollaremos la lectura 9 en $\R^d$, de modo que el caso original $d=1$ y la ampliación multidimensional compartan una misma prueba. La medida es Lebesgue sin normalización.

\section{Existencia, asociatividad e integrales}\label{lectura:v4-09}
\index{convolucion@convolución!euclidiana}
Para funciones medibles definimos $(f*g)(x)=\int_{\R^d}f(y)g(x-y)\,dy$ cuando la integral converge absolutamente. Si converge, el cambio $y=x-z$ prueba $f*g(x)=g*f(x)$. En general la función se define casi en todas partes y representa una clase de equivalencia.
\begin{teorema}\label{v4:conv-Rn}
Si $f,g\in L^1(\R^d)$, la convolución pertenece a $L^1$ y
\[
 \norm{f*g}_1\leq\norm f_1\norm g_1,\qquad
 \int_{\R^d}f*g=\left(\int_{\R^d}f\right)\left(\int_{\R^d}g\right).
\]
La convolución es bilineal, asociativa y conmutativa; hace de $L^1(\R^d)$ un álgebra de Banach.
\end{teorema}
\begin{proof}
Podemos elegir representantes borelianos de las clases. La aplicación $(x,y)\mapsto(y,x-y)$ es un isomorfismo lineal de determinante de valor absoluto uno, por lo que preserva la medida y los nulos. El integrando es medible también para representantes de Lebesgue. Tonelli y la traslación $z=x-y$ dan
\[
 \int_{\R^d}\int_{\R^d}|f(y)g(x-y)|\,dy\,dx
 =\int_{\R^d}|f(y)|\int_{\R^d}|g(z)|\,dz\,dy
 =\norm f_1\norm g_1<\infty.
\]
Fubini asegura existencia casi en todas partes y medibilidad. La desigualdad triangular integral produce la cota de norma; Fubini sin valores absolutos da la identidad de integrales. Cambiar un representante en un nulo no modifica la clase de la convolución, por el mismo cálculo.

Para tres funciones de $L^1$, la integral triple de $|f(z)g(y-z)h(x-y)|$ sobre $(x,y,z)$ es $\norm f_1\norm g_1\norm h_1$. Casi todo $x$ permite, por tanto, intercambiar las integrales en $y,z$. El cambio $y=z+t$ da
\[
 ((f*g)*h)(x)=\int f(z)\int g(t)h(x-z-t)\,dt\,dz=(f*(g*h))(x).
\]
Más generalmente, la misma igualdad vale en un punto fijo $x$ si $(|f|*|g|)*|h|(x)<\infty$, con las integrales internas definidas en las secciones pertinentes. La bilinealidad es inmediata y la completitud ya es la de $L^1$. En la fórmula de integrales de la lectura 9 deben figurar integrales sobre $\R$, no sobre $[0,2\pi]$.
\end{proof}

\section{La desigualdad de Young}
\index{Young!desigualdad de convolucion@desigualdad de convolución}
\begin{lema}[Hölder para varios factores]
Si $p_j\in[1,\infty]$ y $\sum_{j=1}^m1/p_j=1$, entonces $\norm{\prod_j f_j}_1\leq\prod_j\norm{f_j}_{p_j}$.
\end{lema}
\begin{proof}
Los factores de exponente infinito se acotan por su supremo esencial. Para los restantes, después de dividir cada factor por su norma, la desigualdad de Young ponderada da $\prod_j|f_j|\leq\sum_j|f_j|^{p_j}/p_j$. Integrando, el lado derecho vale uno. La desigualdad ponderada se obtiene de la convexidad de la exponencial: $\exp(\sum\theta_j a_j)\leq\sum\theta_j e^{a_j}$, con $\theta_j=1/p_j$ y $a_j=p_j\log|f_j|$; los ceros se tratan por límite. Si alguna norma es cero, el producto se anula casi en todas partes.
\end{proof}
\begin{teorema}[Young]\label{v4:young-Rn}
Sean $p,q,r\in[1,\infty]$ con $1/p+1/q=1+1/r$. Si $f\in L^p(\R^d)$ y $g\in L^q(\R^d)$, entonces
\[
 \norm{f*g}_r\leq\norm f_p\norm g_q.
\]
La convolución existe absolutamente casi en todas partes; si $r=\infty$, puede definirse en todo punto usando representantes adecuados, mediante Hölder.
\end{teorema}
\begin{proof}
Si $r=\infty$, $q=p'$ y Hölder aplicado a $f(y)$ y $g(x-y)$ da la cota para cada $x$. Si $p=1$ y $r=q<\infty$, Hölder para la medida $|f(y)|dy$ da
\[
 |f*g(x)|^q\leq\norm f_1^{q-1}\int|f(y)||g(x-y)|^qdy.
\]
Integrar en $x$ y aplicar Tonelli da la conclusión, incluida existencia casi en todas partes. Para $q=1$ se intercambian los factores. El caso $p=q=r=1$ ya está probado.

Queda $1<p,q\leq r<\infty$. Escribimos
\[
 |f(y)g(x-y)|=
 (|f(y)|^p|g(x-y)|^q)^{1/r}
 |f(y)|^{1-p/r}|g(x-y)|^{1-q/r}.
\]
Los exponentes de Hölder son $r$, $(1/p-1/r)^{-1}$ y $(1/q-1/r)^{-1}$; sus recíprocos suman uno. Un factor con exponente de potencia cero se omite. Como $|f|^p,|g|^q\in L^1$, el teorema anterior asegura que su convolución es finita casi en todas partes. Hölder da en esos puntos
\[
 |f*g(x)|^r\leq\norm f_p^{r-p}\norm g_q^{r-q}(|f|^p*|g|^q)(x).
\]
Integrar y usar la fórmula de la integral de convolución produce $\norm{f*g}_r^r\leq\norm f_p^r\norm g_q^r$. Todos los cambios se han hecho primero sobre integrandos no negativos, por lo que Tonelli justifica también la integrabilidad que se concluye.
\end{proof}
En particular, $T_f:g\mapsto f*g$ es lineal y acotado en $L^p$, con $\norm{T_f}\leq\norm f_1$. La igualdad de norma afirmada en la lectura 9 no vale para todos los exponentes. Por ejemplo, para $p=2$, $f=\mathbf1_{[0,1]}-\mathbf1_{[1,2]}$ tiene norma $L^1$ igual a dos, pero Plancherel dará $\norm{T_f}=\sup|\widehat f|<2$: en cada frecuencia la desigualdad triangular es estricta, pues las fases no pueden ser constantes en ambos intervalos con signos opuestos; la transformada es continua y se anula en infinito, por lo que alcanza su máximo y este es estrictamente menor que dos.

\section{Continuidad y comportamiento en infinito}
\begin{teorema}\label{v4:conv-continuidad}
Si $1\leq p<\infty$, $f\in L^p(\R^d)$ y $g\in L^{p'}(\R^d)$, entonces $f*g$ es acotada y uniformemente continua. Si $1<p<\infty$, además pertenece a $C_0(\R^d)$.
\end{teorema}
\begin{proof}
Hölder da $|f*g(x)|\leq\norm f_p\norm g_{p'}$ para todo $x$. La continuidad de traslaciones en $L^p$, probada en el ejercicio~\ref{ej:v4-l09-1}, implica
\[
 |f*g(x+h)-f*g(x)|\leq\norm g_{p'}\norm{f(\cdot+h)-f}_p\longrightarrow0,
\]
uniformemente en $x$. Si $1<p<\infty$, aproximamos $f,g$ en sus respectivas normas por funciones continuas de soporte compacto $u,v$. Hölder da
\[
 \norm{f*g-u*v}_\infty\leq\norm{f-u}_p\norm g_{p'}+\norm u_p\norm{g-v}_{p'}.
\]
La convolución $u*v$ es continua y se anula fuera de la suma compacta de los soportes. Por ello $f*g$ es límite uniforme de funciones de $C_c$, y pertenece a $C_0$. Para $p=1$, la aproximación de $g\in L^\infty$ por funciones de soporte compacto no está garantizada; el ejemplo $g=1$ explica por qué la conclusión adicional puede fallar.
\end{proof}
\begin{ejemplo}
En $\R$, $\mathbf1_{[0,1]}*\mathbf1_{[0,1]}(x)$ es la longitud de $[0,1]\cap[x-1,x]$. Vale $x$ para $0\leq x\leq1$, $2-x$ para $1\leq x\leq2$, y cero fuera de $[0,2]$. La convolución de dos funciones discontinuas puede así ser continua. En $\R^d$, la convolución de las indicadoras de $[0,1]^d$ es el producto de estas funciones triangulares, por Tonelli.
\end{ejemplo}

\section{Ejercicios: traslaciones, soportes y aproximación}
\begin{ejercicio}[Lectura 9, ejercicio 1]\label{ej:v4-l09-1}
Para $1\leq p<\infty$ y $f\in L^p(\R^d)$, demostrar que $t\mapsto f(\cdot+t)$ y $t\mapsto f(\cdot-t)$ son uniformemente continuas con valores en $L^p$.
\end{ejercicio}

\begin{ejercicio}[Lectura 9, ejercicio 2]\label{ej:v4-l09-2}
Sean $f\in L^1(\R^d)$ y $g$ acotada, con $g(x)\to0$ cuando $|x|\to\infty$. Probar que $f*g(x)\to0$ y mostrar que no puede omitirse una condición de este tipo en un enunciado general.
\end{ejercicio}

\begin{ejercicio}[Lectura 9, ejercicio 3]\label{ej:v4-l09-3}
Si $f,g$ tienen soporte compacto y su convolución existe casi en todas partes, demostrar $\supp(f*g)\subseteq\supp f+\supp g$, entendiendo soporte esencial cuando corresponda.
\end{ejercicio}

\begin{ejercicio}[Lectura 9, ejercicio 4]\label{ej:v4-l09-4}
Sea $(k_j)\subset L^1(\R^d)$, con $k_j\geq0$, $\int k_j=1$ y $\int_{|t|\geq\delta}k_j(t)dt\to0$ para cada $\delta>0$. Probar $f*k_j\to f$ en $L^p$ para $1\leq p<\infty$.
\end{ejercicio}

\begin{ejercicio}[Lectura 9, ejercicio 5: continuidad en los puntos del compacto]\label{ej:v4-l09-5}
Sea $f$ un representante acotado de $L^\infty(\R^d)$ que es continuo, como función definida en $\R^d$, en cada punto de un compacto $K$. Probar que $f*k_j\to f$ uniformemente en $K$. Explicar por qué no basta continuidad de la restricción $f|_K$.
\end{ejercicio}


Las traslaciones y las identidades aproximadas han convertido integrales en límites en norma. Una identidad gaussiana permitirá ahora invertir Fourier mediante cálculos explícitos.
